\documentclass[11pt,a4paper]{article}

% --- Packages ---
\usepackage[margin=.75in]{geometry}
\usepackage[utf8]{inputenc}
\usepackage[T1]{fontenc}
\usepackage{lmodern} % Fix for font not found error
\usepackage{titlesec}
\usepackage{enumitem}
\usepackage{hyperref}
\usepackage{xcolor}
\usepackage{fancyhdr}
\usepackage{comment}
\usepackage{cleveref}
\usepackage[bottom]{footmisc}
\usepackage{graphicx}
\usepackage{longtable}
\usepackage{booktabs}
\usepackage{array}
\usepackage{calc}
\usepackage{etoolbox}

% Standard Pandoc setup
\providecommand{\tightlist}{%
  \setlength{\itemsep}{0pt}\setlength{\parskip}{0pt}}

\crefformat{section}{\S#2#1#3}
\crefformat{subsection}{\S#2#1#3}
\crefformat{subsubsection}{\S#2#1#3}
\crefformat{figure}{#2Figure~#1#3}
\crefformat{footnote}{#2\footnotemark[#1]#3}

\linespread{0.9} % Riduce lo spazio tra le linee

% --- Colors ---
\definecolor{primary}{RGB}{0, 0, 0}
\definecolor{links}{HTML}{0000EE}
\definecolor{blue}{HTML}{26428B}

\hypersetup{
    colorlinks=true,
    linkcolor=links,
    urlcolor=links,
    citecolor=links,
}

% --- Section Formatting ---
\titleformat{\section}
  {\normalfont\Large\bfseries\color{blue}}{\thesection}{1em}{}
\titleformat{\subsection}
{\normalfont\large\bfseries\color{blue}}{\thesubsection}{1em}{}
\titleformat{\subsubsection}
{\normalfont\bfseries\color{blue}}{\thesubsubsection}{1em}{}

% --- Header and Footer ---
\pagestyle{fancy}
\fancyhead[L]{\small Giuseppe Capasso}
\fancyhead[R]{\small intro-to-linux}
\fancyfoot[C]{\thepage}
\renewcommand{\headrulewidth}{0.4pt}

% --- Custom Title Command ---
\newcommand{\proposaltitle}[1]{
    \begin{center}
        \vspace*{0.1cm}
        {\LARGE \bfseries \color{blue} #1} \\[1.5ex]
        {\large \textbf{Giuseppe Capasso}} \\[1ex]
                \small{
            \vspace{0.1cm}
            \textbf{Materia:} intro-to-linux\\
        }
            \end{center}
    \vspace{0.3cm}
    \hrule height 0.5pt
    \vspace{0.5cm}
}

\begin{document}

\proposaltitle{Oppure: sudo ufw allow 22/tcp}

Dispensa Tecnica: Amministrazione Remota, Sicurezza e Systemd Modulo:
Linux System Administration - Lezione 5 Obiettivo: Imparare a gestire un
server Linux ``Headless'' (senza interfaccia grafica) dal proprio PC
fisico, configurare servizi in background e proteggere il sistema con un
firewall.

\begin{enumerate}
\def\labelenumi{\arabic{enumi}.}
\tightlist
\item
  L'Ambiente Server e il ``Port Forwarding'' I veri server non hanno un
  monitor attaccato. Sono macchine inserite in armadi (rack) all'interno
  di datacenter lontani chilometri. Per simularlo, useremo Ubuntu Server
  su VirtualBox e lo gestiremo dal terminale del nostro PC fisico
  (Host). Poiché la nostra Macchina Virtuale (VM) è configurata in
  modalità NAT, si trova dietro un ``router virtuale'' che la isola dal
  nostro PC. Per poterci entrare, dobbiamo creare un ``tunnel'' chiamato
  Port Forwarding (Inoltro delle porte).
\end{enumerate}

1.1 Come configurare il Port Forwarding su VirtualBox Obiettivo: Dire a
VirtualBox che tutto il traffico in arrivo sulla porta 2222 di Windows
deve essere spedito alla porta 22 (SSH) di Linux. 1. Apri VirtualBox,
seleziona la tua VM (anche se è accesa) e clicca su Impostazioni. 2. Vai
nella sezione Rete e assicurati che sia impostato su NAT. 3. Clicca su
Avanzate e poi sul pulsante Inoltro delle porte. 4. Clicca sull'icona +
verde a destra per aggiungere una regola: • Nome: SSH\_Access •
Protocollo: TCP • IP Host: 127.0.0.1 (Indica il tuo computer locale) •
Porta Host: 2222 • IP Guest: (Lascia vuoto) • Porta Guest: 22 (La porta
standard del servizio SSH in Linux) 5. Clicca OK e applica le modifiche.

2. Secure Shell (SSH): Il controllo remoto SSH è il protocollo standard
per amministrare sistemi operativi remoti in modo sicuro. Crea una
connessione crittografata impenetrabile.

2.1 Il Primo Accesso Ora che il Port Forwarding è attivo, riduci a icona
VirtualBox. Apri PowerShell (o il Terminale) su Windows e digita:
PowerShell ssh -p 2222 studente@127.0.0.1

Se ti chiede di accettare un ``fingerprint'', scrivi yes. Poi inserisci
la tua password di Linux.

2.2 Autenticazione a Chiavi (Key-Based Auth) Usare le password è scomodo
e insicuro (soggetto ad attacchi brute-force). I professionisti usano la
Crittografia Asimmetrica. 1. Genera le chiavi sul tuo PC Windows: (Apri
una nuova scheda di PowerShell, non farlo dentro Linux!) PowerShell
ssh-keygen -t ed25519

(Premi sempre Invio per accettare i valori predefiniti). 2. Copia la
chiave pubblica sul Server Linux: Esegui questo comando da PowerShell
per ``spingere'' la tua chiave pubblica dentro la VM: PowerShell cat
\$env:USERPROFILE.ssh\id\_ed25519.pub \textbar{} ssh -p 2222
studente@127.0.0.1 ``mkdir -p \textasciitilde/.ssh \&\& cat
\textgreater\textgreater{} \textasciitilde/.ssh/authorized\_keys''

\begin{enumerate}
\def\labelenumi{\arabic{enumi}.}
\setcounter{enumi}{2}
\tightlist
\item
  Il Test: Chiudi la connessione SSH digitando exit. Ricollegati con ssh
  -p 2222 studente@127.0.0.1. Se non ti chiede la password, hai
  configurato correttamente le chiavi!
\end{enumerate}

2.3 Blindare il Server (Disattivare le Password) Dato che ora entriamo
con la nostra chiave univoca, possiamo dire al server di rifiutare
qualsiasi tentativo di accesso tramite password. Sul terminale Linux
(tramite SSH), edita il file di configurazione: Bash sudo nano
/etc/ssh/sshd\_config

Cerca la riga \#PasswordAuthentication yes, togli il commento (\#) e
modificala in: Plaintext PasswordAuthentication no

Attenzione: questa modifica avrà effetto solo dopo aver riavviato il
servizio SSH (vedi capitolo 4).

\begin{enumerate}
\def\labelenumi{\arabic{enumi}.}
\setcounter{enumi}{2}
\tightlist
\item
  Sicurezza di Base: Il Firewall (UFW) Un server esposto in rete subisce
  continui tentativi di accesso. Dobbiamo usare un Firewall per bloccare
  il traffico non desiderato. Su Ubuntu usiamo UFW (Uncomplicated
  Firewall). REGOLA D'ORO: UFW lavora sulle porte interne di Linux.
  Quindi dovremo aprire la porta 22 (SSH), a prescindere dal fatto che
  su VirtualBox abbiamo usato la 2222.
\end{enumerate}

3.1 La sequenza di attivazione corretta Se attivi il firewall prima di
autorizzare la porta SSH, verrai ``buttato fuori'' dalla tua stessa
connessione remota! Segui questo ordine: 1. Imposta il blocco totale in
entrata (``Default Deny''): Bash sudo ufw default deny incoming sudo ufw
default allow outgoing

\begin{enumerate}
\def\labelenumi{\arabic{enumi}.}
\setcounter{enumi}{1}
\tightlist
\item
  Autorizza i servizi necessari (Whitelist): Bash sudo ufw allow ssh
\end{enumerate}

\section{Oppure: sudo ufw allow
22/tcp}\label{oppure-sudo-ufw-allow-22tcp}

\begin{enumerate}
\def\labelenumi{\arabic{enumi}.}
\setcounter{enumi}{2}
\item
  Attiva il Firewall: Bash sudo ufw enable
\item
  Controlla che tutto sia corretto: Bash sudo ufw status verbose
\end{enumerate}

4. Systemd: Il ``Manager'' dei Servizi I programmi che devono rimanere
sempre accesi in background (come il server web o lo stesso SSH) sono
chiamati Demoni o Servizi. In Linux, questi servizi sono gestiti da un
sistema centrale chiamato Systemd. Il comando per parlare con Systemd è
systemctl.

4.1 Comandi Essenziali Da imparare a memoria per la gestione quotidiana:
• sudo systemctl status {[}servizio{]} ➔ Verifica se il servizio sta
funzionando (cerca la scritta active (running)). • sudo systemctl start
{[}servizio{]} ➔ Accende il servizio. • sudo systemctl stop
{[}servizio{]} ➔ Spegne il servizio. • sudo systemctl restart
{[}servizio{]} ➔ Lo riavvia (necessario dopo aver modificato i file in
/etc/, come abbiamo fatto prima con SSH!). • sudo systemctl enable
{[}servizio{]} ➔ Dice a Linux di accendere questo servizio
automaticamente ogni volta che il server si avvia. Prova subito ad
applicare la modifica fatta nel capitolo 2.3 digitando: sudo systemctl
restart ssh

4.2 Creare un servizio personale da zero Creiamo un piccolo server web
Python e trasformiamolo in un servizio Systemd immortale. 1. Prepara i
file: Bash sudo mkdir /opt/mioservizio echo ``

Sito Web servito da Systemd!

'' \textbar{} sudo tee /opt/mioservizio/index.html

\begin{enumerate}
\def\labelenumi{\arabic{enumi}.}
\setcounter{enumi}{1}
\item
  Crea la ``Unit File'' (Il manuale di istruzioni per Systemd): Bash
  sudo nano /etc/systemd/system/miniweb.service
\item
  Scrivi questa configurazione e salva (CTRL+O, CTRL+X): Ini, TOML
  {[}Unit{]}
\end{enumerate}

Description=Il mio mini server web Python After=network.target

{[}Service{]} Type=simple WorkingDirectory=/opt/mioservizio \# Esegue il
modulo web di Python sulla porta 8000 ExecStart=/usr/bin/python3 -m
http.server 8000 Restart=on-failure

{[}Install{]} WantedBy=multi-user.target

\begin{enumerate}
\def\labelenumi{\arabic{enumi}.}
\setcounter{enumi}{3}
\tightlist
\item
  Fai partire il tuo nuovo Demone: Bash sudo systemctl daemon-reload
\end{enumerate}

\section{Fai leggere a Systemd i nuovi
file}\label{fai-leggere-a-systemd-i-nuovi-file}

sudo systemctl start miniweb

\section{Accendilo}\label{accendilo}

sudo systemctl enable miniweb

\section{Rendilo automatico al
riavvio}\label{rendilo-automatico-al-riavvio}

sudo systemctl status miniweb

\section{Verifica che sia acceso}\label{verifica-che-sia-acceso}

Vuoi vederlo funzionare? Ritorna su VirtualBox (Port Forwarding) e
aggiungi una regola: • Nome: Web\_Access, Porta Host: 8080, Porta Guest:
8000. Poi vai sul terminale Linux e apri la porta nel firewall: sudo ufw
allow 8000. Infine, apri il browser su Windows e vai all'indirizzo
http://127.0.0.1:8080. Se vedi la scritta, sei ufficialmente un
SysAdmin!



\end{document}
